\documentclass[12pt]{ctexart}
\usepackage{a4}
\usepackage{graphicx}
\usepackage{geometry}
\usepackage{amsmath}
\usepackage{mathtools}
\usepackage{amssymb}
\geometry{a4paper,left=2cm,right=2cm,top=2.5cm,bottom=2.5cm}

\title{几何基础作业13}
\author{无名}
\date{\today}

\begin{document}
\maketitle
\section{设直线$l$与双曲线$\Omega$ 交于$A,B$两点,与两条渐近线分别交于$Q,P$
求证$\left\lvert BQ\right\rvert=\left\lvert PA\right\rvert  $}
\begin{figure}[htb]
    \begin{center}
        \includegraphics[width=0.5\linewidth]{1.png}
        \caption{}
    \end{center}
\end{figure}
我们在射影平面$\mathbb{RP^2}$上构造一个射影变换，使得原来的双曲线变为圆，渐近线变为切线，如图
\begin{figure}[htb]
    \begin{center}
        \includegraphics[width=0.5\linewidth]{2.png}
        \caption{}
    \end{center}
\end{figure}
\newpage
那么此时有$(A'B'P'Q')=-1$

根据射影变换交比不变，我们有$(ABPQ)=-1$在原图形内，即
\[\frac{AP}{PB}:\frac{AQ}{QB}=-1\]

记$AB$中点为$M$

\[\frac{AP}{PB}=\frac{AM+MP}{MB-MP}\]
\[\frac{AQ}{QB}=\frac{AM+MQ}{QM+MB}\]

代入到交比中，得到$M$也是$PQ$中点

此时
\[\left\lvert BQ\right\rvert=\left\lvert BM\right\rvert-\left\lvert MQ\right\rvert=\left\lvert AM\right\rvert-\left\lvert MP\right\rvert    =\left\lvert PA\right\rvert \]
\end{document}